\documentclass{article}
\usepackage[paperwidth=16in,paperheight=10.5in,margin=0.5in]{geometry}
\usepackage[T1]{fontenc}
\usepackage{lmodern,amsmath,booktabs,tabularx,array}
\usepackage[table]{xcolor}
\usepackage{hyperref}
\renewcommand{\familydefault}{\sfdefault}
\pagestyle{empty}
\setlength{\parindent}{0pt}
\definecolor{maroon}{HTML}{890B23}
\definecolor{blue}{HTML}{0076BA}
\definecolor{pale}{HTML}{F8EDEF}
\newcolumntype{Y}{>{\raggedright\arraybackslash}X}
\setlength{\tabcolsep}{10pt}
\renewcommand{\arraystretch}{1.35}
\begin{document}
{\centering\fontsize{30}{35}\selectfont\bfseries\color{maroon}
AMP-forming versus ADP-forming synthetases\par}
\vspace{10pt}
{\fontsize{17}{21}\selectfont Two routes to carboxyl activation during ATP-dependent amide ligation}\par
\vspace{18pt}
\fontsize{18}{23}\selectfont
\begin{tabularx}{\textwidth}{@{}>{\bfseries\raggedright\arraybackslash}p{3.7in}YY@{}}
\toprule
\rowcolor{pale}
Feature & \textbf{AMP-forming (adenylation)} & \textbf{ADP-forming (phosphorylation)} \\
\midrule
ATP attack site & $\alpha$-phosphate ($P_\alpha$) & $\gamma$-phosphate ($P_\gamma$) \\
Activated intermediate & Acyl-adenylate (acyl-AMP) & Acyl-phosphate \\
Activation coproduct & Pyrophosphate ($PP_i$) & ADP \\
Aminolysis leaving group & AMP & Inorganic phosphate ($P_i$) \\
Overall ATP products & AMP $+$ $PP_i$ & ADP $+$ $P_i$ \\
Relevant families and functional groups &
\begin{minipage}[t]{\linewidth}\fontsize{16}{21}\selectfont
ANL amide ligases (McbA)\\
GH3 acyl acid amido synthetases\\
NIS synthetases (IucA, IucC, DesD)\\
Asparagine synthetases A and B (AsnA, AsnB)\\
AsnB-related $\beta$-lactam / carbapenam synthetases\\
Other characterized ligases (DdaG)
\end{minipage}
&
\begin{minipage}[t]{\linewidth}\fontsize{16}{21}\selectfont
ATP-grasp peptide / isopeptide ligases (Ddl, DdaF, PylC)\\
ATP-grasp glutathione / polyamine ligases (GshB, GspS, TryS)\\
GS-fold ligases (GS, PuuA, GmaS)\\
Glutamate--cysteine ligases (GCL / GshA)\\
Mur ligases (MurC--MurF, Mpl)\\
FolC / folylpolyglutamate synthetases\\
Cyanophycin synthetase CphA1 (ATP-grasp + Mur-like)
\end{minipage}\\[5pt]
Activation geometry & $\angle(O_{\rm nuc},P_\alpha,O3A)$ & $\angle(O_{\rm nuc},P_\gamma,O3B)$ \\
\bottomrule
\end{tabularx}
\par\vspace{17pt}
{\fontsize{18}{22}\selectfont\bfseries\color{blue}
Assign the activation mechanism before choosing the intermediate and structural proxy.}\par
\vspace{10pt}
{\fontsize{11}{14}\selectfont
Scope: major carboxyl-activating amide-ligase families and functional groups. Entries are not counts of independent folds;
some are related or contain multiple domains. Substrate scope differs among members. Carrier-dependent NRPS, tRNA and
protein-conjugation systems are listed separately in the extended inventory. Metal stoichiometry requires separate evidence.
O3A and O3B are ATP atom names.\\[4pt]
Hollenhorst et al., \textit{Biochemistry} (2009), \href{https://doi.org/10.1021/bi9013165}{doi:10.1021/bi9013165};
Mydy et al., \textit{Biochemistry} (2020), \href{https://pmc.ncbi.nlm.nih.gov/articles/PMC7325057/}{PMCID:PMC7325057};
Kurihara et al., \textit{JBC} (2008), \href{https://pubmed.ncbi.nlm.nih.gov/18495664/}{PMID:18495664}.}
\par\vspace{5pt}
{\fontsize{10}{13}\selectfont
Extended family sources: Petchey et al. (2018), \href{https://doi.org/10.1002/ange.201804592}{doi:10.1002/ange.201804592};
GH3 study (2024), \href{https://pmc.ncbi.nlm.nih.gov/articles/PMC11253546/}{PMC11253546};
AsnB-related ligases, \href{https://pmc.ncbi.nlm.nih.gov/articles/PMC2610506/}{PMC2610506};
GspS, \href{https://pmc.ncbi.nlm.nih.gov/articles/PMC1698887/}{PMC1698887};
GCL, \href{https://pmc.ncbi.nlm.nih.gov/articles/PMC523444/}{PMC523444};
PylC, \href{https://doi.org/10.1016/j.jmb.2012.09.007}{doi:10.1016/j.jmb.2012.09.007};
CphA1, \href{https://doi.org/10.1038/s41589-021-00854-y}{doi:10.1038/s41589-021-00854-y}.
See the extended inventory for row-specific references and substrate-context limitations.}
\end{document}
